\scanpage{63}
Les bitriangles sont en corr. biunivoque avec l'ens des \inserted{couples $(a,u)$ d'un} sommet de $\Delta$, \deleted{en associant à un tel sommet la réunion de} \inserted{et d'un \deleted{\unsure{orb}}} élément $u\in\deleted{\unsure{\mathrm{II}|\Pi}}I(X_\alpha)=I_\alpha$ où $X_\alpha\ni a$ et $I(X_\alpha)=I_\alpha$ est l'ens des trois orbites \unsure{associé|associés|associées}.
À un couple $(a,u)$ correspond le bitriangle formé de $a$, des 4 sommets liés à $a$ dans $X_\alpha$, et des 4 sommets de $S$ qui sont dans l'orbite choisie $u$, et qui sont non liés à $a$.
Pour un $a\in X=X_\alpha$ donné, il y a donc exact\textsuperscript{mt} 3 bitriangles $Y$ qui coupent $X$ suivant la réunion des deux triangles de $X$ de sommet $a$ ; et pour un doublet donné $\{a,b\}$ de $X$ ($a\ne b$, $a$ non lié à $b$) il existe exact\textsuperscript{mt} $\deleted{\unsure{2}}3\deleted{\unsure{6}}$ bitriangles $Y$ coupant $X$ en $\{a,b\}$. [Cela donne

\diagram{$\begin{array}{@{}lr@{}}
\text{\deleted{les trois}}&\\
\text{le bitriangle } X&1\\
\text{les 2 bitriangles } X',X''\ \text{disjoints de } X&2\\
\text{les }3\times9\ \text{bitriangles }Y\text{ coupant }X\text{ suivant une}&27\\
\qquad\text{2-étoile de triangles}&\\
\text{les }\unsure{3|6}\times\deleted{\unsure{10|9}}\unsure{18|9}\ \text{bitriangles }Y\text{ coupant }X\text{ suivant}&54\\
\qquad\text{un doublet}&\deleted{\unsure{90|84}}\\
\text{les }\deleted{\unsure{9\times4}}\inserted{6\times6}\ \text{bitriangles }Y\text{ coupant }X&\\
\qquad\text{suivant un triangle}&36\\ \cline{2-2}
\text{\deleted{\uline{NB} il y a}}&120
\end{array}$}

On avait \deleted{\unsure{oublié}} négligé de détailler le cas des $Y$ engendrés par des triades de type III, plus simple.

% archive sheet number: 62

